% !TeX root = ./../main.tex

\chapter*{Wykaz skrótów i symboli}
\addcontentsline{toc}{chapter}{Wykaz skrótów i symboli}
\agentinfo{
Wykaz uzupełniać na bieżąco wraz z rozdziałami.
Makra symboli znajdują się w config/symbols.tex.
Symbol użyty w tekście musi mieć wiersz w tym wykazie.
}

\section*{Skróty}

\begin{longtable}{@{}p{0.2\textwidth}p{0.73\textwidth}@{}}
\toprule
\textbf{Skrót} & \textbf{Znaczenie} \\
\midrule
\endhead
CPTP & odwzorowanie całkowicie dodatnie i zachowujące ślad (ang. \emph{completely positive trace-preserving}) \\
FRQI & \emph{Flexible Representation of Quantum Images} \\
GAN & generatywna sieć przeciwstawna (ang. \emph{Generative Adversarial Network}) \\
LPIQE & \emph{Local Phase Image Quantum Encoding} \\
MAE & średni błąd bezwzględny (ang. \emph{Mean Absolute Error}) \\
MSE & błąd średniokwadratowy (ang. \emph{Mean Square Error}) \\
NEQR & \emph{Novel Enhanced Quantum Representation} \\
NISQ & \emph{Noisy Intermediate-Scale Quantum} \\
PDF & funkcja rozkładu prawdopodobieństwa (ang. \emph{Probability Distribution Function}) \\
PDO & operator zniekształcenia fazy (ang. \emph{Phase Distortion Operator}) \\
PDU & funkcja odwracania zniekształceń fazy (ang. \emph{Phase Distortion Unraveling}) \\
PSNR & szczytowy stosunek sygnału do szumu (ang. \emph{Peak Signal-to-Noise Ratio}) \\
PTM & macierz transferu Pauliego (ang. \emph{Pauli Transfer Matrix}) \\
QCoSamp & kwantowe próbkowanie szeregiem cosinusów (ang. \emph{Cosine Series Quantum Sampling}) \\
QEC & kwantowa korekcja błędów (ang. \emph{Quantum Error Correction}) \\
QEM & mitygacja błędów kwantowych (ang. \emph{Quantum Error Mitigation}) \\
QML & kwantowe uczenie maszynowe (ang. \emph{Quantum Machine Learning}) \\
QPT & tomografia procesu kwantowego (ang. \emph{Quantum Process Tomography}) \\
QV & objętość kwantowa (ang. \emph{Quantum Volume}) \\
SPAM & błędy przygotowania stanu i pomiaru (ang. \emph{State Preparation and Measurement}) \\
VQC & wariacyjny obwód kwantowy (ang. \emph{Variational Quantum Circuit}) \\
\bottomrule
\end{longtable}

\section*{Symbole}

\begin{longtable}{@{}p{0.2\textwidth}p{0.73\textwidth}@{}}
\toprule
\textbf{Symbol} & \textbf{Znaczenie} \\
\midrule
\endhead
$\hilbert$ & przestrzeń Hilberta \\
$\ket{\psi}$ & wektor stanu układu kwantowego \\
$\rho$ & macierz gęstości \\
$\identity$ & operator identycznościowy \\
$\channel$ & kanał kwantowy \\
$\prob(\ket{x})$ & prawdopodobieństwo zmierzenia stanu bazowego $\ket{x}$ \\
$\reffun(x)$ & funkcja referencyjna o znanych wartościach oczekiwanych \\
$\reffunexp(x)$ & funkcja referencyjna wyznaczona eksperymentalnie \\
$\pdu(x)$ & funkcja PDU: $\pdu(x)=\reffunexp(x)-\reffun(x)$ \\
$H, W$ & wysokość i szerokość obrazu w pikselach \\
$\img$ & obraz jako macierz intensywności pikseli \\
$\lpiqe(\img)$ & operator kodowania LPIQE \\
$\reconst(\img)$ & operator rekonstrukcji obrazu \\
$\liealg{g}$ & algebra Liego grupy $G$ \\
$\genpert$ & addytywne zaburzenie generatora bramki \\
$\dataset$ & zbiór danych \\
$\trainset$, $\valset$, $\testset$ & zbiór treningowy, walidacyjny i testowy \\
$\symloss$ & funkcja straty \\
\bottomrule
\end{longtable}
